\section{Introduction}
\label{submission:sec:intro}

The rising demand for high-throughput multi-task serving has spurred interest in techniques that can dynamically run multiple expert tasks under constrained resource footprints. A prominent approach is \emph{dynamic model merging} or \emph{test-time ensembling}, where task-specific parameter-efficient adapters (such as low-rank adapters or LoRAs~\cite{DudaHart2nd}) are dynamically scaled and blended on a sample-by-sample basis at inference time. Rather than routing queries to isolated full-parameter models, dynamic model merging passes each query through a single shared backbone model, dynamically ensembling the adapter outputs in a single execution pass.

To achieve fast and lightweight ensembling, modern systems employ activation-space routing. For instance, SABLE~\cite{anonymous} projects intermediate representations onto task-specific principal component analysis (PCA) subspaces to extract coordinate signal vectors that drive the ensembling weights. However, stateless projections are highly sensitive to query noise, resulting in high-frequency oscillations of ensembling weights—a phenomenon known as the \emph{routing jitter paradox}. High jitter degrades serving stability, increases prediction entropy, and ruins the mathematical alignment of blended representation spaces across consecutive layers.

To suppress jitter, stateful routing frameworks such as ChemMerge~\cite{anonymous} and PAC-Kinetics~\cite{anonymous} were introduced. These methods model routing trajectories as continuous-time biochemical kinetics or state space models. By integrating a temporal memory state, they smooth out coordinate noise, establishing an elegant temporal consistency across sequentially served queries. 

\begin{figure}[t]
\centering
\includegraphics[width=0.98\columnwidth]{../results/fig1_decoupling_comparison.png}
\caption{Decoupling dynamics of LDS-Kinetics across different network blocks. Early blocks (left) remain highly dynamic and responsive to task transitions, middle blocks (center) balance stability with adaptation, while late blocks (right) learn high temporal inertia to act as stable low-pass decision-making filters.}
\label{fig:decoupling_comparison}
\end{figure}

However, a fundamental limitation remains in all current stateful model merging architectures: they assume \emph{spatial homogeneity} across the depth of the network. That is, they maintain a single, global ensembling coefficient vector $\alpha_t$ at time step $t$ and apply it uniformly across all dynamic ensembling layers. This design overlooks a foundational principle of deep learning: neural representations evolve dramatically along the depth of a network. Early layers typically capture local features and perform transient representational alignment, middle layers perform task integration and semantic abstraction, and late layers refine output logits. Forcing all layers to ensemble experts using a single temporal scale of memory forces a sub-optimal trade-off between adaptation speed and decision stability.

In this paper, we challenge the spatial homogeneity assumption by introducing \textbf{Layer-Decoupled Stateful Kinetics (LDS-Kinetics)}. We treat network depth as an active variable in temporal-spatial ensembling. LDS-Kinetics partitions the dynamic ensembling layers (e.g., Layers 4 to 14 in our 14-layer backbone) into $M$ disjoint blocks. Each block maintains an independent concentration state vector that evolves according to its own block-specific parameters. This allows early layers to adapt rapidly to incoming transitions, while late layers act as stable, high-inertia low-pass filters to prevent logit-level jitter. 

Because decoupling parameters across $M$ blocks dramatically increases the optimization search space, it introduces a significant risk of transductive overfitting on short calibration streams. As empirical researchers, we address this threat by incorporating a unified PAC-Bayesian complexity penalty. Derived from Catoni’s $\beta$-mixing PAC bound~\cite{catoni2007}, this learning-theoretic regularization simultaneously constrains all $M$ sets of decoupled parameters, centering them around safe, default stateless priors.

Adhering to our core philosophy that true scientific progress in machine learning must be built on exhaustive empirical validation, we perform extensive multi-dimensional sweeps to evaluate LDS-Kinetics. Using a 14-layer coordinate sandbox simulator across 5 independent random seeds, we benchmark LDS-Kinetics under both orthogonal and overlapping task manifold environments, and across both homogeneous and heterogeneous sequential workloads. Our contributions are threefold:
\begin{itemize}
    \item We formulate \textbf{LDS-Kinetics}, the first dynamic model merging framework that decouples stateful ensembling dynamics along network depths, generalized to arbitrary block scales $M$.
    \item We integrate a unified learning-theoretic PAC-Bayesian complexity penalty that mathematically bounds and empirically mitigates overfitting in high-dimensional decoupled routing spaces.
    \item We present a comprehensive empirical deconstruction of spatial-temporal routing dynamics. Our sweeps demonstrate that unregularized decoupled routing fails to generalize, and that regularized LDS-Kinetics successfully uncovers depth-dependent kinetics: late blocks learn low decay (high inertia) to stabilize decisions, while early blocks maintain high decay to track rapid workload transitions.
\end{itemize}
